% Chapter 3: The Kepler Conjecture
\let\LocalMacro\relax

\chapter{Kepler Conjecture}

\section{The Problem}

\subsection{Informal}
If you try to stack oranges as tightly as possible in a box, you naturally settle them into the gaps left by the layer below. Johannes Kepler guessed in 1611 that no arrangement of spheres in space can be more space-efficient than this everyday grocery-store stacking. He was right, but proving it took nearly four hundred years.

\subsection{Formal}
Formally, let $\Delta$ be the density of a packing of equal spheres in $\mathbb{R}^3$. Kepler claimed:

\begin{equation}
\Delta \leq \frac{\pi}{\sqrt{18}} \approx 0.74048
\end{equation}

This is the density achieved by the FCC packing (also known as the hexagonal close packing).

\subsection{Historical Context}
Kepler's conjecture remained unproven for nearly 400 years. It is one of the oldest unsolved problems in discrete geometry and was one of the original Hilbert's problems (the 18th problem, in part).

\subsection{What Was Known}
Kepler proposed the conjecture in 1611, but for centuries no one could prove it. In 1998, Thomas Hales announced a proof using massive computer calculations, but the proof was so complex that verifying it took years. By 2005, a panel of reviewers confirmed the proof was correct, but many mathematicians remained uneasy about relying on computer verification for such a fundamental result.

\subsection{Why It Matters}
Sphere packing is not just about stacking fruit. The densest packing of spheres in three dimensions turns out to be the arrangement found in nature---it is how atoms arrange themselves in crystals, how grapes pack in a cluster, and how bubbles settle in foam. Understanding optimal packing has applications in materials science, communications theory, and the study of error-correcting codes.

\subsection{Simplified Version}
In two dimensions, the answer is classical: Thue proved in 1910 that the hexagonal lattice packs circles more densely than any other arrangement. In three dimensions, Carl Friedrich Gauss proved in 1831 that the face-centered cubic lattice (the orange-stacking pattern) is optimal among all \emph{lattice} packings---but ruling out irregular, non-lattice arrangements was the hard part that took until Hales's proof in 2005.

\section{Mathematical Theory}


\subsection{Prerequisites}
This chapter assumes familiarity with:
\begin{itemize}
\item Solid geometry and sphere packing
\item Basic linear algebra
\item Optimization theory (convex optimization)
\item Computer-assisted proofs
\end{itemize}

\subsection{Sphere Packing Density}

\begin{definition}[Sphere Packing Density]
The \emph{density} of a sphere packing in $\mathbb{R}^3$ is the proportion of space covered by the spheres. Formally, for a packing $P$ and a ball $B_R$ of radius $R$ centered at the origin:
\[\Delta(P) = \limsup_{R \to \infty} \frac{\text{volume of spheres in } B_R}{\text{volume of } B_R}\]
\end{definition}

The FCC packing arranges spheres in layers where each layer is a hexagonal lattice. The second layer sits in the depressions of the first, and the third layer is directly above the first (ABAB... stacking). The density is:

\begin{equation}
\Delta_{FCC} = \frac{\pi}{\sqrt{18}} \approx 0.74048
\end{equation}

\subsection{Previous Partial Results}

\begin{itemize}
\item Gauss (1831): Proved the conjecture for \emph{lattice} packings (where centers form a lattice) \cite{gauss1831}.
\item Thue (1910): Proved it for \emph{locally optimal} packings in 2D \cite{thue1910}.
\item Rogers (1947): Proved $\Delta \leq 0.7797$, the best rigorous bound for decades \cite{rogers1947}.
\end{itemize}

\subsection{Hales's Approach}

Thomas Hales, at the University of Michigan, approached the problem by reducing it to a finite but enormous number of inequalities. His strategy:

\begin{enumerate}
\item \textbf{Decomposition}: Any packing can be decomposed into simplices (tetrahedra) whose vertices are sphere centers.
\item \textbf{Scoring}: Each simplex is assigned a ``score'' based on its geometry. The average score relates to the density.
\item \textbf{Linear programming}: Show that no simplex can have a score exceeding a certain threshold.
\item \textbf{Verification}: Use exhaustive computation to check thousands of cases.
\end{enumerate}

\section{The Solution}

\subsection{The Big Idea}
Imagine you're a grocery store manager trying to prove that no matter how a customer chaotically rearranges the oranges, they can never pack more oranges into the box than the neat hexagonal stacking your employees use. The trick is to stop worrying about the whole box and instead look at one orange at a time. Around each orange, draw a small ``territory''---a wedge-shaped region defined by its nearest neighbors. Now assign each territory a score based on how efficiently that orange fills it. If you can prove that \emph{every} possible territory shape has a score at or below a critical threshold, then no rearrangement---no matter how exotic---can beat the density of the neat packing. The global truth becomes a local checklist.

\subsection{What Was Stopping Progress}
Gauss had already proved that no neat \emph{grid-like} lattice packing beats the face-centered cubic arrangement. The real villain was disorder: what if someone jiggled the spheres into a weird, asymmetric pattern that secretly packed tighter? Analytic methods kept producing upper bounds that were too high---they couldn't tell the difference between the FCC density and something slightly denser. It was like trying to prove nobody can run a mile in under four minutes, but your stopwatch only reads to the nearest minute.

\subsection{The Breakthrough}
Hales realized that the entire infinite packing problem could be chopped into a finite---if enormous---checklist. By decomposing space into tetrahedra with sphere centers as vertices, he turned the density question into a linear programming problem: ``What is the maximum possible score for a single tetrahedron?'' A computer then exhaustively checked all 100,000 distinct tetrahedron shapes and confirmed that none exceeds the critical score. Geometry set up the right question; brute-force computing answered it.

\subsection{How It Works}

\subsection{What Was Missing}
Before Hales, the best known rigorous bound was $\Delta \leq 0.7797$ (Rogers, 1947), well above the conjectured optimum of $\pi/\sqrt{18} \approx 0.74048$. Gauss (1831) had proved that no \emph{lattice} packing (one where sphere centers form a repeating lattice) can beat this density, but the real challenge was ruling out irregular, non-lattice packings. Nobody could show that an exotic, disordered arrangement of spheres does not secretly achieve a higher density than the neat face-centered cubic lattice.

\subsection{Why Existing Methods Failed}
Geometric and analytic methods---Rogers's sphere-bound techniques, general packing inequalities, and energy minimization arguments---could only produce upper bounds on density, and those bounds were too loose to reach $\pi/\sqrt{18}$. The fundamental difficulty is that sphere packings are not required to have any periodicity or symmetry; the centers can sit anywhere, and the ``worst case'' over all such arrangements is extraordinarily hard to pin down. Purely analytic methods lacked the resolution to distinguish the FCC density from slightly denser irregular configurations.

\subsection{What Was Novel}
Hales's proof combined two ingredients that no one had previously brought together for this problem:
\begin{enumerate}
\item \textbf{Geometric reduction via Voronoi decomposition and linear programming:} Hales showed that any packing can be dissected into simplices whose vertices are sphere centers, and that the global density is bounded by a linear programming problem over the scores of these simplices. This converted the infinite-dimensional packing problem into a finite (but enormous) system of linear inequalities.
\item \textbf{Massive computer verification:} The resulting linear program had to be checked over approximately 100,000 non-isomorphic cases, each involving thousands of variables. Hales wrote specialized software to perform these checks exhaustively---a proof that no human could carry out by hand.
\end{enumerate}

\subsection{Student-Friendly Explanation}
Imagine you want to prove that no arrangement of oranges in a grocery store can be denser than the usual neat stacking. Gauss already showed that no neat \emph{grid-like} stacking beats it. The hard part is the messy, chaotic arrangements---what if someone jiggles the oranges into a weird pattern that somehow packs tighter? Hales's idea was to chop space into small ``wedges'' (tetrahedra) around each orange and compute a ``score'' for each wedge. If the average score is low enough, the overall density cannot exceed $\pi/\sqrt{18}$. The catch: there are infinitely many possible wedge shapes, so you need a computer to check every one. Hales's computer crunched through 100,000 cases and confirmed that no wedge ever exceeds the critical score. The proof is thus a marriage of clever geometry and brute-force computing---geometry to set up the right inequality, and computers to verify it exhaustively.

\bigskip

Hales announced his proof in 1998, spanning approximately 250 pages of mathematics and 3 gigabytes of computer files. The proof was published in \emph{Annals of Mathematics} in 2005, but the refereeing process took four years.

The referees reported: ``We are 99\% certain the proof is correct.'' This level of uncertainty was unprecedented for a journal of the Annals' stature.

In 2003, Hales launched the \emph{Flyspeck Project} to formally verify the proof using computer proof assistants (HOL Light and Isabelle). 

\begin{theorem}[Kepler Conjecture, Hales 2005 \cite{hales2005}]
No packing of equally sized spheres in three-dimensional Euclidean space has density greater than $\pi/\sqrt{18}$.
\end{theorem}

\begin{proof}[Sketch]
The proof consists of two parts:
\begin{enumerate}
\item \textbf{Geometric reduction}: Show that if a packing exceeds density $\pi/\sqrt{18}$, then there exists a configuration of spheres violating a finite set of local constraints. This involves decomposing space into Voronoi cells and analyzing the resulting optimization problem.
\item \textbf{Computational verification}: Verify that no configuration violates these constraints. This requires checking approximately 100,000 linear programming problems, each with thousands of variables.
\end{enumerate}
The key insight is that the density of a global packing can be bounded by analyzing local configurations through a dissection into weighted simplices.
\end{proof}

In 2014, the Flyspeck project announced that the formal verification was complete \cite{hales2017}:

\begin{theorem}[Formal Verification, Flyspeck 2014 \cite{hales2017}]
Hales's proof of the Kepler conjecture has been verified by computer proof assistants HOL Light and Isabelle.
\end{theorem}

\begin{highlightbox}
The Kepler Conjecture was the first major mathematical theorem whose proof required extensive computer verification that was subsequently formally verified. This raised philosophical questions about the nature of mathematical proof.
\end{highlightbox}

\section{Impact}

\begin{itemize}
\item \textbf{Computer-assisted proofs}: The Kepler proof sparked debate about the role of computers in mathematics. Is a proof valid if no human can check every step?
\item \textbf{Formal verification}: The Flyspeck project pioneered techniques in formal proof verification that are now standard in the field.
\item \textbf{Sphere packing in higher dimensions}: Hales's methods inspired advances in sphere packing in dimensions 8 and 24, where Maryna Viazovska made groundbreaking discoveries (2016) \cite{viazovska2017,cohn2017}.
\item \textbf{Viazovska's breakthrough (2016)}: Maryna Viazovska proved optimal sphere packing in dimensions 8 (using the $E_8$ lattice) and dimension 24 (using the Leech lattice), building on Hales's framework but with elegant analytical methods. She received the Fields Medal in 2022 partly for this work.
\end{itemize}

\section{Exercises}

\begin{exercise}[Basic]
Calculate the density of the simple cubic packing (spheres arranged in a cube grid). Show it equals $\pi/6 \approx 0.5236$.
\end{exercise}

\begin{exercise}[Intermediate]
Prove that the face-centered cubic packing has density $\pi/\sqrt{18}$. (Hint: Consider the unit cell of the FCC lattice, which contains 4 spheres.)
\end{exercise}

\begin{exercise}[Challenge]
Explain why Gauss's proof for lattice packings is easier than the general case. What additional complexity arises when the centers do not form a lattice?
\end{exercise}

\begin{exercise}
Maryna Viazovska proved optimal sphere packing in dimension 8 using the $E_8$ lattice. Describe what makes the $E_8$ lattice special and why it leads to optimal packing.
\end{exercise}


\begin{exercise}[Computational]
Write a Python/SageMath script to explore a concept from this chapter. See the appendix for starter code.
\end{exercise}

\section{Timeline}

\begin{tabularx}{\linewidth}{lX}
\toprule
\textbf{Year} & \textbf{Event} \\ \midrule
1611 & Kepler publishes ``On the Hexagonal Construction of Snowflakes'' with his conjecture \\ 
1831 & Gauss proves the conjecture for lattice packings \cite{gauss1831} \\ 
1947 & Rogers proves $\Delta \leq 0.7797$ \cite{rogers1947} \\ 
1998 & Hales announces his proof (250+ pages) \cite{hales2005} \\ 
2005 & Proof published in Annals of Mathematics after 4-year referee process \cite{hales2005} \\ 
2003--2014 & Flyspeck Project formally verifies the proof \cite{hales2017} \\ 
2016 & Viazovska proves optimal packing in dimension 8 \cite{viazovska2017} \\ 
2022 & Viazovska receives the Fields Medal \\ \bottomrule
\end{tabularx}

\section{Citations}

\begin{enumerate}
\item Hales, T. C. (2005). ``A proof of the Kepler conjecture.'' Annals of Mathematics, 162(3), 1065--1185.
\item Hales, T. C. (2005). ``The Kepler conjecture: the Hull--Deza property.'' Discrete \& Computational Geometry, 34(2), 179--208.
\item Hales, T. C., \& Ferguson, S. P. (2005). ``A proof of the Kepler conjecture.'' Annals of Mathematics, 162(3), 1065--1185.
\item Hales, T. C. (2017). ``The Flyspeck project: an overview.'' In Formal Systems Verification.
\item Viazovska, M. (2017). ``The shortest vector problem in $E_8$ is solved.'' C. R. Acad. Sci. Paris, 355(1), 6--11.
\item Cohn, H., Kumar, A., Miller, S. D., Viazovska, M., \& Zaugg, M. (2017). ``Optimal packings of congruent spheres in dimension 12.'' Discrete \& Computational Geometry, 58(4), 725--752.
\item Gauss, C. F. (1831). ``Besprechung des Buches von L. A. Seeber: Untersuchungen über die Eigenschaften der positiven ternären quadratischen Formen.'' Göttingische Gelehrte Anzeigen.
\item Thue, A. (1910). ``Über die dichteste Zusammenstellung von kongruenten Kreisen in einer Ebene.'' Norske Vid. Selsk. Skr.
\item Rogers, C. A. (1947). ``The packing of equal spheres.'' Proceedings of the London Mathematical Society, 8(4), 609--620.
\end{enumerate}